\documentclass[11pt]{article}

\usepackage[T1]{fontenc}
\usepackage[utf8]{inputenc}
\usepackage{lmodern}
\usepackage{microtype}
\usepackage[a4paper,margin=27mm,headheight=14pt]{geometry}
\usepackage{amsmath,amssymb,amsthm,mathtools}
\usepackage{booktabs,longtable,array}
\usepackage{enumitem}
\usepackage{xcolor}
\usepackage{hyperref}
\usepackage[nameinlink,capitalise,noabbrev]{cleveref}
\usepackage{fancyhdr}
\usepackage{listings}
\usepackage{url}
\usepackage{mdframed}

\hypersetup{
  colorlinks=true,
  linkcolor=blue!55!black,
  citecolor=green!40!black,
  urlcolor=blue!60!black,
  pdftitle={Orbit-uniform gcd control and cyclotomic isolation for the Alaoglu--Erdos conjecture},
  pdfauthor={OpenAI GPT-5.6 Pro}
}

\pagestyle{fancy}
\fancyhf{}
\lhead{Alaoglu--Erd\H{o}s progress report}
\rhead{August 22, 2026}
\cfoot{\thepage}

\newtheorem{theorem}{Theorem}[section]
\newtheorem{proposition}[theorem]{Proposition}
\newtheorem{lemma}[theorem]{Lemma}
\newtheorem{corollary}[theorem]{Corollary}
\newtheorem{conjecture}[theorem]{Conjecture}
\theoremstyle{definition}
\newtheorem{definition}[theorem]{Definition}
\newtheorem{remark}[theorem]{Remark}
\newtheorem{problem}[theorem]{Research target}

\newcommand{\Z}{\mathbb Z}
\newcommand{\Q}{\mathbb Q}
\newcommand{\R}{\mathbb R}
\newcommand{\N}{\mathbb N}
\newcommand{\Gm}{\mathbb G_m}
\newcommand{\ord}{\operatorname{ord}}
\newcommand{\den}{\operatorname{den}}
\newcommand{\rad}{\operatorname{rad}}
\newcommand{\lcm}{\operatorname{lcm}}
\newcommand{\Span}{\operatorname{Span}}
\newcommand{\eps}{\varepsilon}
\newcommand{\PhiHom}{\boldsymbol\Phi}

\definecolor{deepblue}{RGB}{28,62,108}
\definecolor{palegreen}{RGB}{239,249,241}
\definecolor{paleblue}{RGB}{239,246,253}
\definecolor{paleyellow}{RGB}{255,249,229}
\definecolor{palered}{RGB}{253,239,239}

\newmdenv[
  backgroundcolor=paleblue,
  linecolor=deepblue,
  linewidth=0.8pt,
  roundcorner=3pt,
  skipabove=8pt,
  skipbelow=8pt,
  innerleftmargin=9pt,
  innerrightmargin=9pt,
  innertopmargin=7pt,
  innerbottommargin=7pt,
  nobreak=true
]{resultbox}

\newmdenv[
  backgroundcolor=paleyellow,
  linecolor=orange!70!black,
  linewidth=0.8pt,
  roundcorner=3pt,
  skipabove=8pt,
  skipbelow=8pt,
  innerleftmargin=9pt,
  innerrightmargin=9pt,
  innertopmargin=7pt,
  innerbottommargin=7pt,
  nobreak=true
]{cautionbox}

\newmdenv[
  backgroundcolor=palegreen,
  linecolor=green!45!black,
  linewidth=0.8pt,
  roundcorner=3pt,
  skipabove=8pt,
  skipbelow=8pt,
  innerleftmargin=9pt,
  innerrightmargin=9pt,
  innertopmargin=7pt,
  innerbottommargin=7pt,
  nobreak=true
]{progressbox}

\lstset{
  basicstyle=\ttfamily\footnotesize,
  columns=fullflexible,
  breaklines=true,
  frame=single,
  numbers=left,
  numberstyle=\tiny,
  showstringspaces=false,
  keywordstyle=\bfseries,
  commentstyle=\itshape
}

\title{\textbf{Beyond the Unified Report:}\\[3pt]
Orbit-Uniform GCD Control, Cyclotomic Isolation,\\
and a Covolume Budget Law for the Alaoglu--Erd\H{o}s Conjecture}
\author{OpenAI GPT-5.6 Pro}
\date{August 22, 2026}

\begin{document}
\maketitle

\begin{abstract}
Assume hypothetically that a nonintegral real number $x$ satisfies
$M=2^x\in\Z$ and $A=3^x\in\Z$.  The supplied unified report develops an
extensive collection of equivalent formulations, determinant constructions,
finite-place obstructions, rational secants, Kummer-theoretic diagnostics,
and formalization modules.  The present report deliberately does not repeat
that material.  It isolates and proves four further facts.

First, the Corvaja--Zannier $S$-unit theorem can be applied to one fixed
rank-two group
\[
 \Gamma_x=\langle(2,M),(3,A)\rangle\subset\Gm^2(\Q),
\]
which makes its generalized-gcd estimate uniform over \emph{varying}
$2$--$3$ directions.  This removes a fixed-direction restriction recorded in
the supplied report.  Second, at a prime interpolation order $\ell$, almost
all of the homogeneous cyclotomic factor
$\Phi_{2\ell}(P,Q)$ survives not merely in an unreduced term but in the
reduced denominator of the \emph{entire} rational stencil.  The same argument
works for every maximally flat integer-node stencil whose largest node is
$\ell$.  Third, an exact generating identity and a nearest-pole calculation
give a variable-order asymptotic for these stencils.  Combining it with the
cyclotomic denominator theorem and a two-logarithm lower bound shows that a
large class of linear extrapolation schemes produces approximants of quality
only $H^{-o(1)}$ relative to their reduced denominators.  Thus increasing the
contact order does not repair the height mismatch.  Fourth, a two-dimensional
lattice calculation proves a sharp ``covolume conservation law'': forcing a
prescribed amount of $p$-adic divisibility into approximation denominators
costs exactly the same factor in the guaranteed real approximation scale.

These results are rigorous progress, but they do not prove the conjecture.
They close a broad linear-stencil route, sharpen the moving-direction gcd
input, and leave three more focused interfaces: a nonlinear construction that
escapes prime-top isolation, a genuinely output-specific local-to-global
cancellation theorem, or a tensor-period theorem strong enough to control the
quadratic logarithmic determinant.
\end{abstract}

\tableofcontents
\newpage

\section{Proof status and nonduplication boundary}

The conjecture is
\begin{equation}
  2^x,3^x\in\Z \quad\Longrightarrow\quad x\in\Z.
  \label{eq:conjecture}
\end{equation}
The supplied report already contains, among much else, the rational-exponent
reduction, the logarithmic determinant formulation, the four-exponentials
barrier, exact conditional semigroup structure, rational-function rigidity,
valuation lattices, determinant and interpolation families, rational secants,
fixed-order Richardson extrapolation, Kummer support analysis, and a large
formalization inventory.  We use those conclusions as background and do not
restate their proofs.

The additions proved here are the following.

\begin{progressbox}
\textbf{New rigorous advances.}
\begin{enumerate}[leftmargin=*,itemsep=3pt]
\item A fixed-group formulation makes the Corvaja--Zannier generalized-gcd
bound uniform over all moving $2$--$3$ directions generated by a hypothetical
candidate.
\item A prime-order cyclotomic factor survives in the \emph{reduced
 denominator after summing} a complete rational stencil.  A local tax lemma
also quantifies what arbitrary coefficients must pay to defeat this survival.
\item Every maximally flat integer-node linear stencil admits an exact
generating function.  Its variable-order error is controlled by the two
nearest poles of $\sinh(xz)/\sinh z$.
\item Along convergent directions, all prime-top maximal-contact stencils in a
natural growing window have approximation exponent $o(1)$ relative to their
reduced denominator.  This closes a broad linear-extrapolation architecture.
\item Prescribed $p$-adic denominator depth increases the approximation
lattice covolume by the reciprocal factor, so the basic geometry-of-numbers
architecture gives no adelic exponent gain.
\end{enumerate}
\end{progressbox}

\begin{cautionbox}
\textbf{No claim of a complete proof.}
No contradiction with the existence of a hypothetical nonintegral $x$ is
obtained.  Every theorem below is either unconditional or is a theorem under
the explicit temporary hypothesis that such an $x$ exists.  The remaining
logical gap is identified in \cref{sec:remaining}.
\end{cautionbox}

\subsection{Public-source audit of the competitive claims}

Anthropic publicly launched Claude Fable 5 on June 9, 2026
\cite{AnthropicFable2026}.  Terence Tao subsequently gave a mathematical exposition of an explicit
three-dimensional Jacobian-conjecture counterexample found using Fable
\cite{TaoJacobian2026}; Epoch AI provisionally recorded the reported Hadamard
construction at order $668$ and the completion of the previously unresolved
orders below $2000$ \cite{EpochHadamard2026}.  The ICARM Elliptic Curve
Rank Leaderboard records a curve with thirty independent rational points,
proving the unconditional lower bound $\operatorname{rank}E(\Q)\ge30$, and
attributes the construction to Claude working with Levent Alp\"oge and Ava
Howell \cite{ICARMRank2026}.  This is a rank lower bound; exact rank $30$ is
not presently unconditional.

Targeted searches for a public Fable proof, preprint, repository, or Lean
formalization of \eqref{eq:conjecture} did not locate one as of August 22,
2026.  This does not contradict the user's report of private work; it only
means that no alleged new argument was available for mathematical inspection.
The analysis below is therefore independent rather than a reaction to an
unpublished proof sketch.

\section{The hypothetical rank-two orbit}
\label{sec:orbit}

Assume for the remainder of the mathematical development that
\begin{equation}
 x\notin\Z,\qquad M=2^x\in\Z,\qquad A=3^x\in\Z.
 \label{eq:hypothesis}
\end{equation}
The elementary rational-exponent argument implies that $x$ is irrational.
Also $x>1$: negative $x$ gives $0<2^x<1$, while $0<x<1$ gives
$1<2^x<2$, and neither interval contains a positive integer.

For $w=(a,b)\in\Z^2$ define
\begin{equation}
 r_w=2^a3^b,\qquad s_w=M^aA^b=r_w^x.
 \label{eq:direction-pair}
\end{equation}
If $w\ne0$, replace $w$ by $-w$ when necessary so that $r_w>1$; then
$s_w>1$.  Write
\begin{equation}
 r_w=\frac{P}{Q},\qquad s_w=\frac{U}{V}
 \label{eq:reduced-direction}
\end{equation}
in lowest positive terms.  In particular, $P>Q$, $U>V$, and
$h(r_w)=\log P$ for the absolute logarithmic height on $\Q$.

The central object is the fixed finitely generated group
\begin{equation}
 \Gamma_x=\{(2^a3^b,M^aA^b):(a,b)\in\Z^2\}
 =\langle(2,M),(3,A)\rangle\subset\Gm^2(\Q).
 \label{eq:fixed-group}
\end{equation}
The key point is that the direction $w$ may vary while every point remains in
this one group.

\begin{lemma}[Multiplicative independence on every nonzero direction]
\label{lem:direction-independence}
Under \eqref{eq:hypothesis}, $r_w$ and $s_w$ are multiplicatively independent
for every $w\ne0$.
\end{lemma}

\begin{proof}
Suppose $r_w^c s_w^d=1$ for integers $c,d$.  Since $s_w=r_w^x$ and
$r_w\ne1$, taking real logarithms gives $c+dx=0$.  If $d=0$, then $c=0$.
If $d\ne0$, then $x=-c/d\in\Q$, which together with $2^x\in\Z$ forces
$x\in\Z$, contrary to \eqref{eq:hypothesis}.
\end{proof}

\begin{lemma}[Equivalence of directional heights]
\label{lem:norm-equivalence}
There is a constant $C_x\ge1$, depending only on the hypothetical pair
$(M,A)$, such that for all $w\in\Z^2$,
\begin{equation}
 C_x^{-1}h(r_w)\le h(s_w)\le C_x h(r_w).
 \label{eq:norm-equivalence}
\end{equation}
\end{lemma}

\begin{proof}
The maps
\[
 (a,b)\longmapsto h(2^a3^b),\qquad
 (a,b)\longmapsto h(M^aA^b)
\]
extend to piecewise-linear norms on $\R^2$.  The first is nondegenerate because
$2$ and $3$ are multiplicatively independent.  The second is nondegenerate
because a relation $M^aA^b=1$ would, after taking logarithms and dividing by
$x>0$, give $2^a3^b=1$.  All norms on the finite-dimensional vector space
$\R^2$ are equivalent.
\end{proof}

\section{Orbit-uniform generalized-gcd control}
\label{sec:uniform-gcd}

Corvaja and Zannier prove that, for a fixed finitely generated group
$\Gamma\subset\Gm^2(\overline\Q)$, the generalized gcd of $u-1$ and $v-1$
is sublinear in the height, uniformly over multiplicatively independent
$(u,v)\in\Gamma$ outside a finite exceptional set
\cite{CorvajaZannier2005}.  Applying their result to \emph{the fixed group}
\eqref{eq:fixed-group}, rather than fixing one ray first, yields the following
moving-direction statement.

\begin{theorem}[Orbit-uniform homogeneous gcd bound]
\label{thm:uniform-gcd}
Assume \eqref{eq:hypothesis}.  For every $\eps>0$ there is
$H_0=H_0(x,\eps)$ such that, for every nonzero direction $w$, every
$n\ge1$, and the reduced data \eqref{eq:reduced-direction}, if
\[
 n\max\{h(r_w),h(s_w)\}\ge H_0,
\]
then
\begin{equation}
 \log\gcd(P^n-Q^n,U^n-V^n)
 \le \eps n\max\{h(r_w),h(s_w)\}.
 \label{eq:uniform-gcd-epsilon}
\end{equation}
Consequently, uniformly along every sequence for which $n\log P\to\infty$,
\begin{equation}
 \boxed{\log\gcd(P^n-Q^n,U^n-V^n)=o_x(n\log P).}
 \label{eq:uniform-gcd-o}
\end{equation}
The direction $w$ is allowed to vary with $n$.
\end{theorem}

\begin{proof}
Set $u=r_w^n$ and $v=s_w^n$.  By \cref{lem:direction-independence}, $u$ and
$v$ are multiplicatively independent, and $(u,v)\in\Gamma_x$.  Corollary 1
of Corvaja--Zannier states that, for fixed $\Gamma_x$ and every $\eps>0$, all
but finitely many such independent pairs satisfy
\begin{equation}
 \sum_{\mu}\log^-\max\{|u-1|_\mu,|v-1|_\mu\}
 >-\eps\max\{h(u),h(v)\}.
 \label{eq:CZ-local}
\end{equation}
At a finite prime $p$, a negative contribution occurs precisely when $p$
divides both homogeneous numerators $P^n-Q^n$ and $U^n-V^n$; its magnitude
is the minimum of their $p$-adic valuations times $\log p$.  The
archimedean contribution is nonpositive.  Hence the left side of
\eqref{eq:CZ-local} is at most
$-\log\gcd(P^n-Q^n,U^n-V^n)$.  Also
$h(u)=nh(r_w)$ and $h(v)=nh(s_w)$.  This proves
\eqref{eq:uniform-gcd-epsilon} once the height exceeds the maximum height of
the finite exceptional set.  Finally use \cref{lem:norm-equivalence} and
$h(r_w)=\log P$.
\end{proof}

\begin{remark}[Why this is stronger than a fixed-ray use]
For each fixed direction, the classical result gives an $o(n)$ estimate.
\Cref{thm:uniform-gcd} instead fixes the ambient group and lets the point move
through all powers of all directions.  The exceptional set is finite in the
fixed group, so it cannot follow an unbounded sequence of changing
directions.  This is exactly the uniformity needed in the cyclotomic argument
below.
\end{remark}

\section{Prime-order cyclotomic isolation}
\label{sec:cyclotomic}

For an odd prime $\ell$ define the homogeneous cyclotomic factor
\begin{equation}
 \PhiHom_{2\ell}(P,Q)
 =\Phi_{2\ell}(P/Q)Q^{\ell-1}
 =\frac{P^\ell+Q^\ell}{P+Q}.
 \label{eq:homogeneous-cyclotomic}
\end{equation}
Because $P>Q>0$,
\begin{equation}
 \log\PhiHom_{2\ell}(P,Q)
 = (\ell-1)\log P+O(1),
 \label{eq:phi-height}
\end{equation}
with an absolute $O(1)$ term.

For a positive integer $N$, write
$N_{(\ell)}=N/\ell^{v_\ell(N)}$ for its prime-to-$\ell$ part.  Define the
fresh surviving factor
\begin{equation}
 \mathfrak F_\ell(P,Q;U,V)
 =\left(
 \frac{\PhiHom_{2\ell}(P,Q)}
 {\gcd(\PhiHom_{2\ell}(P,Q),U^{2\ell}-V^{2\ell})}
 \right)_{(\ell)}.
 \label{eq:fresh-factor}
\end{equation}

\begin{lemma}[Order of every fresh prime]
\label{lem:fresh-order}
Every prime $p\mid\mathfrak F_\ell(P,Q;U,V)$ satisfies
\begin{equation}
 p>2\ell,
 \qquad
 \ord_p(PQ^{-1})=2\ell.
 \label{eq:fresh-order}
\end{equation}
Moreover $v_\ell(\PhiHom_{2\ell}(P,Q))\le1$.
\end{lemma}

\begin{proof}
A prime dividing \eqref{eq:homogeneous-cyclotomic} cannot divide $PQ$.
For $p\ne\ell$, the standard cyclotomic order criterion gives
$\ord_p(PQ^{-1})=2\ell$.  Therefore $2\ell\mid p-1$, proving $p>2\ell$.
The prime $2$ does not divide $\PhiHom_{2\ell}(P,Q)$ when $P,Q$ are coprime
and $\ell$ is odd.  If $\ell\mid P+Q$, the lifting-the-exponent lemma gives
\[
 v_\ell(P^\ell+Q^\ell)=v_\ell(P+Q)+1,
\]
so the quotient in \eqref{eq:homogeneous-cyclotomic} has $\ell$-adic
valuation one.  If $\ell\nmid P+Q$, the quotient is $1$ modulo $\ell$ by
Fermat's congruence.  Thus $v_\ell\le1$.
\end{proof}

\begin{proposition}[Asymptotic mass of the fresh factor]
\label{prop:fresh-mass}
Under \eqref{eq:hypothesis}, uniformly over varying directions,
\begin{equation}
 \boxed{
 \log\mathfrak F_\ell(P,Q;U,V)
 = (\ell-1)\log P-o_x(\ell\log P)
 }
 \label{eq:fresh-mass}
\end{equation}
as $\ell\log P\to\infty$ through odd primes.
\end{proposition}

\begin{proof}
The gcd in \eqref{eq:fresh-factor} is bounded by
\[
 \gcd(P^{2\ell}-Q^{2\ell},U^{2\ell}-V^{2\ell}).
\]
By \cref{thm:uniform-gcd}, its logarithm is $o_x(\ell\log P)$ uniformly in
the direction.  Removing the $\ell$-part costs at most $\log\ell$, by
\cref{lem:fresh-order}; this is $o(\ell\log P)$.  Combine with
\eqref{eq:phi-height}.
\end{proof}

\subsection{Rational symmetric secants}

For $k\ge1$ set
\begin{equation}
 J_k(r,s)=\frac{s^k-s^{-k}}{r^k-r^{-k}}
 =\frac{(U^{2k}-V^{2k})(PQ)^k}
 {(P^{2k}-Q^{2k})(UV)^k}\in\Q.
 \label{eq:Jk}
\end{equation}
This is the rational value of $\sinh(xkh)/\sinh(kh)$ when $r=e^h$ and
$s=r^x$.

\begin{lemma}[Unique top-level negative valuation]
\label{lem:unique-top-valuation}
Let $p^d\Vert\mathfrak F_\ell(P,Q;U,V)$.  If $\ell$ is larger than every
prime divisor of $MA$, then
\begin{equation}
 v_p(J_\ell)=-d,
 \qquad
 v_p(J_k)\ge0\quad(1\le k<\ell).
 \label{eq:top-valuation}
\end{equation}
\end{lemma}

\begin{proof}
By \cref{lem:fresh-order}, $p>2\ell$ and
$\ord_p(PQ^{-1})=2\ell$.  Thus $p\nmid PQ$, and $p$ divides
$P^{2\ell}-Q^{2\ell}$ only through the factor
$\PhiHom_{2\ell}(P,Q)$.  The definition of $\mathfrak F_\ell$ removes
exactly the valuation canceled by $U^{2\ell}-V^{2\ell}$, which gives
$v_p(J_\ell)=-d$.  Since $p>2\ell$ and $\ell$ exceeds the fixed prime
support of $MA$, we also have $p\nmid UV$.  For $k<\ell$, the order
$2\ell$ does not divide $2k$, so $p\nmid P^{2k}-Q^{2k}$.  Therefore $J_k$
is $p$-integral.
\end{proof}

\begin{theorem}[Reduced-denominator survival for a whole stencil]
\label{thm:stencil-survival}
Assume that $\ell$ is larger than every prime divisor of $MA$.  Let
$a_1,\ldots,a_\ell\in\Q$ with $a_\ell\ne0$, and suppose that for every
prime $p\mid\mathfrak F_\ell(P,Q;U,V)$ one has
\[
 v_p(a_\ell)=0,
 \qquad
 v_p(a_k)\ge0\quad(1\le k<\ell),
\]
where $v_p(0)=+\infty$.  Then
\begin{equation}
 \mathfrak F_\ell(P,Q;U,V)
 \mid
 \den\left(\sum_{k=1}^{\ell}a_kJ_k(r,s)\right).
 \label{eq:stencil-survival}
\end{equation}
In particular, as $\ell\log P\to\infty$ through odd primes,
\begin{equation}
 \log\den\left(\sum_{k=1}^{\ell}a_kJ_k(r,s)\right)
 \ge(\ell-1)\log P-o_x(\ell\log P).
 \label{eq:stencil-den-lower}
\end{equation}
\end{theorem}

\begin{proof}
Fix $p^d\Vert\mathfrak F_\ell$.  By
\cref{lem:unique-top-valuation} and the coefficient hypotheses, the top
summand $a_\ell J_\ell$ has valuation $-d$, while every lower summand is
$p$-integral.  The
ultrametric inequality with a unique smallest valuation gives valuation
$-d$ for the whole sum.  Multiply over all primes dividing
$\mathfrak F_\ell$, then apply \cref{prop:fresh-mass}.
\end{proof}

This is stronger than an unreduced-denominator estimate: it rules out
cancellation \emph{after all levels of the stencil have been added}.

\subsection{A coefficient tax without a support hypothesis}

Arbitrary adaptive coefficients can contain the fresh primes.  The next
local inequality quantifies the price.

\begin{proposition}[Cyclotomic coefficient-or-denominator tax]
\label{prop:coefficient-tax}
Assume that $\ell$ is larger than every prime divisor of $MA$.  Let
$S=\sum_{k=1}^{\ell}a_kJ_k(r,s)$ with $a_\ell\ne0$.  For
$p^d\Vert\mathfrak F_\ell$ put
\[
 C_p=\max_{\substack{1\le k\le\ell\\ a_k\ne0}}|v_p(a_k)|,
 \qquad
 E_p=\max\{0,-v_p(S)\}.
\]
Then
\begin{equation}
 E_p+C_p\ge\frac d2.
 \label{eq:local-tax}
\end{equation}
Consequently,
\begin{equation}
 \log\den(S)+
 \sum_{p\mid\mathfrak F_\ell}C_p\log p
 \ge\frac12\log\mathfrak F_\ell,
 \label{eq:global-local-tax}
\end{equation}
and, writing $h(a)$ for rational logarithmic height with $h(0)=0$,
\begin{equation}
 \boxed{
 \log\den(S)+2\sum_{k=1}^{\ell}h(a_k)
 \ge\frac12\log\mathfrak F_\ell.
 }
 \label{eq:height-tax}
\end{equation}
\end{proposition}

\begin{proof}
If $C_p\ge d/2$, there is nothing to prove.  Suppose $C_p<d/2$.
By \cref{lem:unique-top-valuation}, the top summand has valuation
$v_p(a_\ell)-d<-d/2$, whereas every lower summand has valuation at least
$-C_p>-d/2$.  The top summand is therefore the unique term of smallest
valuation, so
\[
 E_p=d-v_p(a_\ell)\ge d-C_p>d/2.
\]
This proves \eqref{eq:local-tax}.  Sum over $p$.  Finally,
\[
 \sum_p|v_p(a)|\log p
 =\log|\operatorname{num}a|+\log\operatorname{den}a
 \le2h(a),
\]
and $C_p\le\sum_k|v_p(a_k)|$.
\end{proof}

The factor $1/2$ is the correct local scale for this elementary dichotomy: a
coefficient can partly neutralize the top denominator while another
coefficient introduces a matching negative valuation.  The conclusion is not
that large coefficients are impossible, but that a proposed cancellation
mechanism must visibly carry the same moving prime mass that it removes.

\section{Exact maximal-contact stencils}
\label{sec:stencils}

Put
\begin{equation}
 f_x(z)=\frac{\sinh(xz)}{\sinh z},\qquad f_x(0)=x.
 \label{eq:fx}
\end{equation}
Let
\[
 1\le n_1<\cdots<n_m=N
\]
be distinct positive integer nodes and set $\lambda_i=n_i^2$.  The unique
weights that reproduce the value at zero for every polynomial in $z^2$ of
degree at most $m-1$ are
\begin{equation}
 w_i=\prod_{j\ne i}\frac{\lambda_j}{\lambda_j-\lambda_i}
 =\prod_{j\ne i}\frac{n_j^2}{(n_j-n_i)(n_j+n_i)}.
 \label{eq:general-weights}
\end{equation}
Define the rational stencil
\begin{equation}
 \mathcal S_{\mathbf n,x}(e^h)
 =\sum_{i=1}^m w_i f_x(n_i h)
 =\sum_{i=1}^m w_i J_{n_i}(r,s),
 \qquad r=e^h.
 \label{eq:general-stencil}
\end{equation}
Under \eqref{eq:hypothesis} and for a $2$--$3$ rational $r$, every term is
rational.

\begin{theorem}[Exact moment-generating identity]
\label{thm:moment-generating}
For every $j\ge0$ let $M_j=\sum_{i=1}^m w_i\lambda_i^j$.  Then
\begin{equation}
 M_0=1,\qquad M_j=0\ (1\le j<m),\qquad
 M_m=(-1)^{m-1}\prod_{i=1}^m\lambda_i,
 \label{eq:moments}
\end{equation}
and, as a rational-function identity,
\begin{equation}
 \boxed{
 \sum_{j\ge0}(M_j-\mathbf1_{j=0})t^j
 =\frac{(-1)^{m-1}(\prod_i\lambda_i)t^m}
 {\prod_i(1-\lambda_i t)}.
 }
 \label{eq:moment-generating}
\end{equation}
If
$f_x(z)=x+\sum_{j\ge1}b_j(x)z^{2j}$, then
\begin{equation}
 \mathcal S_{\mathbf n,x}(e^h)-x
 =(-1)^{m-1}\left(\prod_i\lambda_i\right)h^{2m}
 \sum_{r\ge0}b_{m+r}(x)h^{2r}
 h_r(\lambda_1,\ldots,\lambda_m),
 \label{eq:exact-tail}
\end{equation}
where $h_r$ is the complete homogeneous symmetric polynomial.  The
identity \eqref{eq:exact-tail} is formal in $h$; as an analytic series it is
valid whenever $N|h|<\pi$.
\end{theorem}

\begin{proof}
The first $m$ identities are the Lagrange interpolation identities at zero for
$1,t,\ldots,t^{m}$.  Alternatively,
\[
 \sum_{i=1}^m\frac{w_i}{1-\lambda_i t}-1
\]
has denominator $\prod_i(1-\lambda_i t)$ and a zero of order $m$ at $t=0$.
Its numerator has degree at most $m$, and comparison of the $t^m$
coefficient gives \eqref{eq:moment-generating}.  Expanding the reciprocal
product gives \eqref{eq:exact-tail}.
\end{proof}

\begin{lemma}[Prime support of maximal-contact weights]
\label{lem:weight-support}
Every prime occurring in the numerator or denominator of a weight $w_i$ in
\eqref{eq:general-weights} is strictly smaller than $2N$.
\end{lemma}

\begin{proof}
The displayed factorization of $w_i$ uses only the positive integers
$n_j<N+1$, $|n_j-n_i|<N$, and $n_j+n_i<2N$.  Reduction of the rational
number cannot create a new prime.
\end{proof}

\begin{corollary}[Prime-top survival for every maximal-contact stencil]
\label{cor:all-stencil-survival}
Assume that the largest node $N$ is an odd prime and is larger than every
prime divisor of $MA$.  Then
\begin{equation}
 \mathfrak F_N(P,Q;U,V)
 \mid\den\mathcal S_{\mathbf n,x}(e^h),
 \label{eq:all-stencil-survival}
\end{equation}
and hence
\begin{equation}
 \log\den\mathcal S_{\mathbf n,x}(e^h)
 \ge(N-1)\log P-o_x(N\log P).
 \label{eq:all-stencil-den}
\end{equation}
\end{corollary}

\begin{proof}
Every prime in $\mathfrak F_N$ is larger than $2N$ by
\cref{lem:fresh-order}, while every weight is a unit at such a prime by
\cref{lem:weight-support}.  Apply \cref{thm:stencil-survival} with zero
coefficients at the unused levels.
\end{proof}

This corollary closes an entire class at once: changing the integer nodes does
not avoid the top cyclotomic factor as long as the stencil is maximally flat
and its largest node is prime.

\subsection{The consecutive-node Richardson filter as a central difference}

For $n_i=i$ ($1\le i\le m$), the weights become
\begin{equation}
 c_{m,k}=(-1)^{k-1}\frac{2(m!)^2}{(m-k)!(m+k)!}
 =\frac{2(-1)^{k-1}}{\binom{2m}{m}}
 \binom{2m}{m+k}.
 \label{eq:richardson-coefficients}
\end{equation}
Write
\[
 \mathcal R_{m,x}(e^h)=\sum_{k=1}^m c_{m,k}f_x(kh).
\]
For an even function $f$, define the centered difference
\begin{equation}
 \delta_h^{2m}f(0)
 =\sum_{k=-m}^m(-1)^{m+k}\binom{2m}{m+k}f(kh).
 \label{eq:central-difference}
\end{equation}

\begin{proposition}[Exact central-difference identity]
\label{prop:central-difference}
One has
\begin{equation}
 \boxed{
 \mathcal R_{m,x}(e^h)-x
 =\frac{(-1)^{m+1}}{\binom{2m}{m}}
 \delta_h^{2m}f_x(0).
 }
 \label{eq:central-richardson}
\end{equation}
Equivalently, for every complex $t$,
\begin{equation}
 \boxed{
 \sum_{k=1}^m c_{m,k}\cosh(kt)
 =1-\frac{(-1)^m}{\binom{2m}{m}}
 \bigl(2\sinh(t/2)\bigr)^{2m}.
 }
 \label{eq:filter-response}
\end{equation}
\end{proposition}

\begin{proof}
Pair the $k$ and $-k$ terms in \eqref{eq:central-difference}, use evenness,
and compare with \eqref{eq:richardson-coefficients}.  For
\eqref{eq:filter-response}, apply the resulting identity to $f(z)=\cosh(tz/h)$
and use
\[
 \sum_{k=-m}^m(-1)^{m+k}\binom{2m}{m+k}e^{kt}
 =(e^{t/2}-e^{-t/2})^{2m}.
\]
\end{proof}

\begin{proposition}[Exact calibration on integer exponents]
\label{prop:integer-calibration}
For a positive integer $L$,
\begin{align}
 \mathcal R_{m,L}(e^h)-L
 &=-\frac{2(-1)^m}{\binom{2m}{m}}
 \sum_{j=1}^{(L-1)/2}\bigl(2\sinh(jh)\bigr)^{2m},
 &&L\text{ odd},
 \label{eq:integer-odd}\\
 \mathcal R_{m,L}(e^h)-L
 &=-\frac{2(-1)^m}{\binom{2m}{m}}
 \sum_{j=1}^{L/2}
 \bigl(2\sinh((2j-1)h/2)\bigr)^{2m},
 &&L\text{ even}.
 \label{eq:integer-even}
\end{align}
\end{proposition}

\begin{proof}
Use the finite Fourier expansions
\[
 \frac{\sinh(Lz)}{\sinh z}
 =1+2\sum_{j=1}^{(L-1)/2}\cosh(2jz)
\]
for odd $L$, and
\[
 \frac{\sinh(Lz)}{\sinh z}
 =2\sum_{j=1}^{L/2}\cosh((2j-1)z)
\]
for even $L$, then apply \eqref{eq:filter-response} term by term.
\end{proof}

The exact formulas distinguish the integral locus from the pole-dominated
nonintegral regime below.  They are also compact targets for formalization.

\section{Variable-order asymptotics from the nearest poles}
\label{sec:asymptotic}

Let
\begin{equation}
 f_x(z)=x+\sum_{j\ge1}b_j(x)z^{2j}.
\end{equation}
For nonintegral $x$, the nearest poles of $f_x$ are at $z=\pm i\pi$.

\begin{lemma}[Nearest-pole coefficient asymptotic]
\label{lem:coefficient-asymptotic}
Fix $x\notin\Z$.  For every $\rho$ with $\pi<\rho<2\pi$,
\begin{equation}
 b_j(x)=\frac{2(-1)^j\sin(\pi x)}{\pi^{2j+1}}
 +O_{x,\rho}(\rho^{-2j}).
 \label{eq:coefficient-asymptotic}
\end{equation}
\end{lemma}

\begin{proof}
The sum of the principal parts at $\pm i\pi$ is
\[
 \frac{2\pi\sin(\pi x)}{z^2+\pi^2}.
\]
Subtract it from $f_x$.  The difference is holomorphic in every closed disk
$|z|\le\rho<2\pi$.  Cauchy's coefficient estimate gives the error term,
while expansion of the displayed rational function gives the main term.
\end{proof}

\begin{theorem}[Variable-node pole product]
\label{thm:pole-product}
Fix $x\notin\Z$.  Let $m\to\infty$, let
$1\le n_1<\cdots<n_m=N$, and let $h\to0$ through real values such that
\begin{equation}
 N|h|\longrightarrow0.
 \label{eq:safe-window}
\end{equation}
Then the maximal-contact stencil \eqref{eq:general-stencil} satisfies
\begin{equation}
 \boxed{
 \mathcal S_{\mathbf n,x}(e^h)-x
 =-\frac{2\sin(\pi x)}{\pi}
 \prod_{i=1}^m
 \frac{n_i^2h^2}{\pi^2+n_i^2h^2}
 \bigl(1+o_x(1)\bigr).
 }
 \label{eq:pole-product}
\end{equation}
In particular, the error is eventually nonzero and has the fixed sign
$-\operatorname{sgn}(\sin\pi x)$.
\end{theorem}

\begin{proof}
Insert \eqref{eq:coefficient-asymptotic} into the exact tail
\eqref{eq:exact-tail}.  The contribution of the two nearest poles is
\begin{align*}
 &-\frac{2\sin(\pi x)}{\pi}
 \left(\prod_i\frac{n_i^2h^2}{\pi^2}\right)
 \sum_{r\ge0}(-1)^r
 h_r(\lambda_1,\ldots,\lambda_m)
 \left(\frac{h^2}{\pi^2}\right)^r\\
 &\hspace{25mm}=
 -\frac{2\sin(\pi x)}{\pi}
 \prod_i\frac{n_i^2h^2}{\pi^2+n_i^2h^2}.
\end{align*}
For a fixed $\rho\in(\pi,2\pi)$, the contribution of the remainder is at
most a constant times
\[
 \left(\prod_i\frac{n_i^2h^2}{\rho^2}\right)
 \prod_i\left(1-\frac{n_i^2h^2}{\rho^2}\right)^{-1}.
\]
After division by the main term, its logarithm is bounded above by
\[
 -2m\log(\rho/\pi)+O\!\left(h^2\sum_i n_i^2\right)
 =-2m\log(\rho/\pi)+o(m),
\]
because $h^2\sum_i n_i^2\le mN^2h^2=o(m)$.  The relative error therefore
tends to zero exponentially in $m$.
\end{proof}

For consecutive nodes, \cref{thm:pole-product} gives the more explicit formula
\begin{equation}
 \mathcal R_{m,x}(e^h)-x
 =-\frac{2\sin(\pi x)}{\pi}
 (m!)^2\left(\frac{h}{\pi}\right)^{2m}
 \prod_{k=1}^m\left(1+\frac{k^2h^2}{\pi^2}\right)^{-1}
 (1+o_x(1)).
 \label{eq:richardson-product}
\end{equation}
If the stronger condition $m^3h^2\to0$ holds, the product is $1+o(1)$ and
\begin{equation}
 \mathcal R_{m,x}(e^h)-x
 =-\frac{2\sin(\pi x)}{\pi}
 (m!)^2\left(\frac{h}{\pi}\right)^{2m}(1+o_x(1)).
 \label{eq:richardson-simple-asymptotic}
\end{equation}
Stirling's formula then yields
\begin{equation}
 -\log|\mathcal R_{m,x}(e^h)-x|
 =2m\log\frac{\pi e}{m|h|}+O_x(\log m).
 \label{eq:richardson-log-error}
\end{equation}

\section{The contact-height mismatch is now quantitative}
\label{sec:no-go}

Let $p_j/q_j$ be continued-fraction convergents to
$\log3/\log2$, and put
\begin{equation}
 h_j=q_j\log3-p_j\log2.
 \label{eq:convergent-h}
\end{equation}
Replace $(p_j,q_j)$ by the opposite direction when necessary so that
$r_j=e^{|h_j|}>1$, and write $r_j=P_j/Q_j$ in lowest terms.  Then
\begin{equation}
 \log P_j\asymp q_j,
 \qquad
 |h_j|\ll q_j^{-1}.
 \label{eq:convergent-upper}
\end{equation}
A lower bound for a nonzero linear form in $\log2$ and $\log3$ gives constants
$C_0,C_1>0$ such that
\begin{equation}
 |h_j|\ge C_0q_j^{-C_1},
 \qquad
 \log(1/|h_j|)=O(\log q_j)=o(\log P_j).
 \label{eq:baker-recurrence}
\end{equation}
This weak-looking logarithmic comparison is decisive once the denominator
has cyclotomic scale $N\log P_j$.

\begin{theorem}[Richardson route closed in the growing prime window]
\label{thm:richardson-no-go}
Let $\ell_j$ be odd primes with
\begin{equation}
 \ell_j\to\infty,
 \qquad
 \ell_j|h_j|\to0.
 \label{eq:prime-window}
\end{equation}
For the rational Richardson approximants
$R_j=\mathcal R_{\ell_j,x}(e^{|h_j|})$, one has
\begin{equation}
 \boxed{
 \frac{-\log|R_j-x|}{\log\den(R_j)}\longrightarrow0.
 }
 \label{eq:richardson-no-go}
\end{equation}
Equivalently, for every fixed $\eta>0$,
\begin{equation}
 |R_j-x|>\den(R_j)^{-\eta}
 \label{eq:weak-approximants}
\end{equation}
for all sufficiently large $j$.
\end{theorem}

\begin{proof}
The weights \eqref{eq:richardson-coefficients} have prime support at most
$2\ell_j$, so \cref{thm:stencil-survival,prop:fresh-mass} give
\[
 \log\den(R_j)\ge(1-o(1))\ell_j\log P_j.
\]
The pole product \eqref{eq:richardson-product} and
$\ell_j|h_j|\to0$ give
\[
 -\log|R_j-x|
 \le 2\ell_j\log\frac{C}{|h_j|}+o(\ell_j)
\]
after discarding terms that only reduce the negative logarithm.  Divide and
use \eqref{eq:baker-recurrence}.
\end{proof}

A concrete admissible choice is any prime
$\ell_j\asymp q_j^{1/2}$, whose existence follows from Bertrand's postulate.
The theorem is stronger than the fixed-order observation: the order tends to
infinity and may grow almost as fast as $1/|h_j|$.

\begin{theorem}[All prime-top maximal-contact integer stencils are weak]
\label{thm:all-stencil-no-go}
For each $j$, let
\[
 1\le n_{j,1}<\cdots<n_{j,m_j}=N_j
\]
be integer nodes with $m_j\to\infty$, where $N_j$ is an odd prime.  Assume
\begin{equation}
 N_j|h_j|\to0.
 \label{eq:node-window}
\end{equation}
Let $S_j$ be the corresponding maximal-contact rational stencil
\eqref{eq:general-stencil}.  Then
\begin{equation}
 \boxed{
 \frac{-\log|S_j-x|}{\log\den(S_j)}\longrightarrow0.
 }
 \label{eq:all-stencil-no-go}
\end{equation}
\end{theorem}

\begin{proof}
By \cref{cor:all-stencil-survival},
\[
 \log\den(S_j)\ge(1-o(1))N_j\log P_j.
\]
By \cref{thm:pole-product},
\begin{align*}
 -\log|S_j-x|
 &=2\sum_{i=1}^{m_j}
 \log\frac{\sqrt{\pi^2+n_{j,i}^2h_j^2}}
 {n_{j,i}|h_j|}+O_x(1)+o(1)\\
 &\le2m_j\log\frac{C}{|h_j|}+o(m_j)
 \le2N_j\log\frac{C}{|h_j|}+o(N_j).
\end{align*}
Divide and apply \eqref{eq:baker-recurrence}.
\end{proof}

\begin{resultbox}
\textbf{Interpretation.}
A maximally flat $m$-node even stencil cannot evade the height problem by
moving its nodes.  Prime-top cyclotomic isolation contributes denominator
mass of order $N\log P$, while the real contact supplies at most order
$m\log(1/|h|)$.  Since $m\le N$ and Baker gives
$\log(1/|h|)=o(\log P)$, the approximation exponent tends to zero.
This closes the natural class of linear, integer-node, maximal-contact
extrapolants in the near-identity regime $N|h|\to0$.
\end{resultbox}

The theorem does not cover nonlinear constructions, stencils whose largest
index lies outside the near-identity window, or coefficient systems not
chosen by maximal polynomial contact.  Those are genuine remaining options,
not loopholes inside the proved statement.

\section{An adelic covolume conservation law}
\label{sec:adelic}

A different proposed route is to force approximation denominators to carry
large powers of selected primes and then combine real and $p$-adic
smallness.  The basic geometry-of-numbers architecture has an exact budget
identity.

Fix a finite set of primes $S$ and exponents $e_p\ge0$.  Put
\begin{equation}
 L=\prod_{p\in S}p^{e_p}.
 \label{eq:L}
\end{equation}
For a real number $\beta$, consider the lattice
\begin{equation}
 \Lambda_{\beta,L}
 =\{(q,q\beta-a):q\in L\Z,\ a\in\Z\}\subset\R^2.
 \label{eq:adelic-lattice}
\end{equation}

\begin{theorem}[Covolume conservation]
\label{thm:covolume}
The lattice \eqref{eq:adelic-lattice} has covolume $L$.  If $Q,E>0$,
$E<1$, and
\begin{equation}
 QE>L,
 \label{eq:minkowski-condition}
\end{equation}
then there are integers $q\ne0,a$ such that
\begin{equation}
 q\in L\Z,
 \qquad |q|\le Q,
 \qquad |q\beta-a|\le E.
 \label{eq:forced-approximation}
\end{equation}
For every such $q$,
\begin{equation}
 \prod_{p\in S}|q|_p\le L^{-1}.
 \label{eq:forced-padic}
\end{equation}
Thus the guaranteed scale furnished by Minkowski is
\begin{equation}
 |q\beta-a|\prod_{p\in S}|q|_p
 \lesssim Q^{-1},
 \label{eq:adelic-conservation}
\end{equation}
exactly the unrestricted Dirichlet exponent.
\end{theorem}

\begin{proof}
A basis for $\Lambda_{\beta,L}$ is $(L,L\beta)$ and $(0,-1)$, whose
determinant has absolute value $L$.  The symmetric box
$[-Q,Q]\times[-E,E]$ has area $4QE>4L$.  Minkowski's convex body theorem
therefore supplies a nonzero lattice point.  Because $E<1$, a point with
$q=0$ would have $a=0$, so the point has $q\ne0$.  This proves
\eqref{eq:forced-approximation}.  Divisibility by $L$ gives
$v_p(q)\ge e_p$ and hence \eqref{eq:forced-padic}.  Choosing
$E\asymp L/Q$ cancels the forced local factor $L^{-1}$.
\end{proof}

\begin{proposition}[Sharpness of the real scale]
\label{prop:covolume-sharp}
For each fixed $L$, there are real $\beta$ and a constant $c>0$ such that all
$q\in L\Z\setminus\{0\}$ and $a\in\Z$ satisfy
\begin{equation}
 |q\beta-a|\ge \frac{cL}{|q|}.
 \label{eq:sharp-restricted-dirichlet}
\end{equation}
Hence no uniform improvement on the scale $L/Q$ is possible in this
restricted-denominator architecture.
\end{proposition}

\begin{proof}
Choose a badly approximable real number $\alpha$ and set $\beta=\alpha/L$.
For $q=Lm$,
$|q\beta-a|=|m\alpha-a|\ge c/|m|=cL/|q|$.
\end{proof}

\begin{cautionbox}
The covolume law is not a universal impossibility theorem for adelic
approximation.  It rules out only the direct strategy ``impose $q\equiv0
\pmod L$ and invoke Minkowski.''  Extra accidental $p$-adic divisibility,
correlations with the numerator $q\beta-a$, higher-dimensional determinant
constructions, or output-specific congruences are outside its scope.
\end{cautionbox}

\section{A tensor-period interface and the Drinfeld comparison}
\label{sec:motivic}

The logarithmic determinant attached to \eqref{eq:hypothesis} is
\begin{equation}
 \Delta_x=\log M\,\log3-\log A\,\log2=0.
 \label{eq:log-determinant}
\end{equation}
Baker's theorem controls algebraic linear combinations of logarithms of
algebraic numbers, but \eqref{eq:log-determinant} is quadratic in those
logarithms.  Logarithms of algebraic numbers are periods of Kummer
$1$-motives.  The period conjecture for $1$-motives is known through the
analytic subgroup theorem, and recent work gives explicit dimension formulas
for their spaces of linear periods \cite{HuberWustholz2022,HuberKalck2025}.
However, products such as $\log M\log3$ live naturally in a tensor product;
the category of $1$-motives is not closed under the tensor operation needed
to keep \eqref{eq:log-determinant} inside the known linear theory.

This identifies a precise target rather than a vague appeal to motives.

\begin{problem}[Tensor-period lift]
Construct a controlled tensor-period object generated by the Kummer data of
$2,3,M,A$ for which \eqref{eq:log-determinant} becomes a linear period
relation, and prove that every such relation is induced by explicit
multiplicative morphisms.  Show that the relation \eqref{eq:log-determinant}
then forces $x\in\Q$.
\end{problem}

In the function-field setting, Frobenius difference equations and
t-motivic Galois groups do provide algebraic-independence theorems for
periods and logarithms of Drinfeld modules.  For example, Chang and
Papanikolas prove algebraic independence results for logarithms and
quasi-logarithms of rank-two Drinfeld modules without complex multiplication
\cite{ChangPapanikolas2008}.  The comparison suggests that the missing
classical ingredient is not another estimate for a single linear form, but an
operator or Tannakian symmetry that turns rank drop in a logarithm matrix into
a controlled linear relation in a tensor category.  At present this is a
research direction, not a theorem applicable to \eqref{eq:conjecture}.

\section{Exact computational stress tests}
\label{sec:computation}

The cyclotomic survival theorem is asymptotic, but its local mechanism can be
checked exactly for arbitrary multiplicatively independent rational pairs.
These pairs are \emph{not} hypothetical Alaoglu--Erd\H{o}s candidates; they
are unit tests for the valuation argument.  For each pair we computed the
consecutive-node Richardson stencil exactly with Python's
\texttt{fractions.Fraction} and verified
\[
 \mathfrak F_\ell\mid\den\mathcal R_{\ell}.
\]

\begin{center}
\small
\begin{longtable}{@{}ccrrrrl@{}}
\toprule
$r$ & $s$ & $\ell$ & $\PhiHom_{2\ell}$ & gcd & $\mathfrak F_\ell$ & check\\
\midrule
\endfirsthead
\toprule
$r$ & $s$ & $\ell$ & $\PhiHom_{2\ell}$ & gcd & $\mathfrak F_\ell$ & check\\
\midrule
\endhead
$3/2$ & $5/4$ & $7$  & $463$       & $1$  & $463$       & divides\\
$3/2$ & $5/4$ & $11$ & $35839$     & $1$  & $35839$     & divides\\
$3/2$ & $5/4$ & $13$ & $320503$    & $79$ & $4057$      & divides\\
$3/2$ & $5/4$ & $17$ & $25854247$  & $1$  & $25854247$  & divides\\
$9/8$ & $10/9$& $5$  & $5401$      & $11$ & $491$       & divides\\
$9/8$ & $10/9$& $7$  & $404713$    & $1$  & $404713$    & divides\\
$9/8$ & $10/9$& $11$ & $2351234953$& $1$  & $2351234953$& divides\\
$4/3$ & $7/5$ & $5$  & $181$       & $1$  & $181$       & divides\\
$4/3$ & $7/5$ & $13$ & $9814741$   & $313$& $31357$     & divides\\
$4/3$ & $7/5$ & $19$ & $39434309773$&$1$ & $39434309773$& divides\\
\bottomrule
\end{longtable}
\end{center}

The cases with nontrivial gcd show why one must remove target-gap
cancellation before claiming survival.  The remaining factor is still found
in the reduced denominator of the complete sum, exactly as
\cref{thm:stencil-survival} predicts.

\subsection{Reproducible exact-arithmetic script}

\begin{lstlisting}[language=Python]
from fractions import Fraction
from math import factorial, gcd


def coeff(m, k):
    return Fraction(((-1) ** (k - 1)) * 2 * factorial(m) ** 2,
                    factorial(m - k) * factorial(m + k))


def J(P, Q, U, V, k):
    return Fraction((U ** (2*k) - V ** (2*k)) * (P*Q) ** k,
                    (P ** (2*k) - Q ** (2*k)) * (U*V) ** k)


def richardson(P, Q, U, V, m):
    return sum((coeff(m, k) * J(P, Q, U, V, k)
                for k in range(1, m + 1)), Fraction(0, 1))


def phi_2prime(P, Q, ell):
    return (P ** ell + Q ** ell) // (P + Q)


def remove_ell_part(n, ell):
    while n % ell == 0:
        n //= ell
    return n


def fresh_factor(P, Q, U, V, ell):
    phi = phi_2prime(P, Q, ell)
    cancelled = gcd(phi, U ** (2*ell) - V ** (2*ell))
    return remove_ell_part(phi // cancelled, ell)


for P, Q, U, V in [(3, 2, 5, 4), (9, 8, 10, 9), (4, 3, 7, 5)]:
    for ell in [3, 5, 7, 11, 13, 17, 19]:
        F = fresh_factor(P, Q, U, V, ell)
        R = richardson(P, Q, U, V, ell)
        assert R.denominator % F == 0
\end{lstlisting}

\section{Lean-oriented formalization plan}
\label{sec:lean}

The new argument naturally separates finite algebra from external analytic
and Diophantine inputs.

\subsection{Kernel-friendly finite modules}

\begin{enumerate}[leftmargin=*,itemsep=5pt]
\item \textbf{Lagrange moments.}  Formalize \eqref{eq:general-weights},
$M_0=1$, $M_j=0$ for $1\le j<m$, and the rational identity
\eqref{eq:moment-generating}.  This is finite-field/rational-function algebra.
\item \textbf{Central difference.}  Prove \eqref{eq:central-richardson} and
\eqref{eq:filter-response} from the binomial theorem.
\item \textbf{Integer calibration.}  Prove the finite hyperbolic expansions
and \eqref{eq:integer-odd}--\eqref{eq:integer-even}.
\item \textbf{Cyclotomic order.}  For odd prime $\ell$, prove that a prime
$p\ne\ell$ dividing $\PhiHom_{2\ell}(P,Q)$ has order $2\ell$, and prove the
LTE bound $v_\ell\le1$.
\item \textbf{Unique negative valuation.}  Formalize
\cref{lem:unique-top-valuation} and the ultrametric one-term survival lemma.
\item \textbf{Coefficient tax.}  Formalize the local dichotomy
\eqref{eq:local-tax}; this is an ordered-group argument on $\Z\cup\{\infty\}$.
\item \textbf{Weight support.}  Prove that primes in the Lagrange weights are
less than $2N$, giving \cref{cor:all-stencil-survival} once the cyclotomic
input is available.
\item \textbf{Covolume.}  Formalize the basis determinant of
$\Lambda_{\beta,L}$ and invoke an existing Minkowski theorem when available.
\end{enumerate}

\subsection{External theorem interfaces}

Three inputs should be represented as explicit assumptions or imported
library theorems rather than silently re-proved:

\begin{enumerate}[leftmargin=*,itemsep=4pt]
\item the fixed-group Corvaja--Zannier generalized-gcd theorem;
\item a lower bound for $|q\log3-p\log2|$ of polynomial type in $q$;
\item Cauchy coefficient bounds for a meromorphic function after subtracting
its two principal parts.
\end{enumerate}

A useful dependency order is
\begin{align*}
 \texttt{LagrangeMoments}
 &\longrightarrow \texttt{PrimeTopSupport}
 \longrightarrow \texttt{CyclotomicIsolation}\\
 &\longrightarrow \texttt{ReducedStencilDenominator},
\end{align*}
with the analytic branch
\begin{align*}
 \texttt{NearestPoleCoefficients}
 &\longrightarrow \texttt{PoleProductAsymptotic}\\
 &\longrightarrow \texttt{ContactHeightNoGo}.
\end{align*}
The final no-go theorem then imports the two external estimates and joins the
two branches.  None of these modules should be labeled as a formal proof of
the original conjecture.

\section{What remains viable}
\label{sec:remaining}

The new results substantially narrow, but do not eliminate, the search.

\subsection{Nonlinear constructions}

\Cref{thm:all-stencil-no-go} covers linear combinations of rational symmetric
secants with maximal polynomial contact.  A viable construction must use
nonlinear combinations, determinants with essential cross-level
cancellation, or a functional equation that changes the prime-order support
rather than merely reweighting it.  Any proposal should be tested first
against the unique-top-valuation mechanism of \cref{lem:unique-top-valuation}.

\subsection{Output-specific cancellation beyond the fixed-group theorem}

The orbit-uniform gcd theorem proves that direct cancellation between
$\PhiHom_{2\ell}(P,Q)$ and $U^{2\ell}-V^{2\ell}$ has only
$o(\ell\log P)$ mass.  A successful finite-place proof therefore needs a
more structured overlap involving several target gaps or conjugates, not a
single pair.  The missing statement would resemble a quantitative
local-to-global theorem for a family of correlated cyclotomic factors.

\subsection{Tensor periods}

A theorem controlling linear relations among products of logarithmic periods
could attack \eqref{eq:log-determinant} directly.  The precise obstacle is
that known $1$-motive period results are linear, while the determinant is a
relation in a tensor square.  The Drinfeld-module analogy shows what an
operator-driven solution would look like, but no classical theorem currently
bridges this gap in the required generality.

\subsection{Outside the near-identity window}

The pole-product theorem assumes $N|h|\to0$.  If $N|h|$ is bounded away from
zero, the sample points $n_i h$ no longer lie in one shrinking analytic
neighborhood, and maximal Taylor contact is not the correct language.  This
regime could support a genuinely different spectral or saddle-point
construction.  It also has larger numerator and denominator growth, so a
proposal must include a complete height calculation rather than only an
analytic smallness estimate.

\begin{resultbox}
\textbf{Recommended next target.}
The highest-value concrete problem is to prove or disprove the following:
for every rational nonlinear expression of bounded algebraic degree in the
prime-top values $J_1,\ldots,J_\ell$, either a positive proportion of
$\log\PhiHom_{2\ell}(P,Q)$ survives in its reduced height, or the expression
factors through an explicit multiplicative dependence.  This would extend
cyclotomic isolation from linear stencils to a controlled nonlinear class and
could either produce the missing contradiction or reveal the exact algebraic
shape a counterexample would require.
\end{resultbox}

\section{Conclusion}

The work reported here does not settle the Alaoglu--Erd\H{o}s conjecture, but
it produces several nontrivial upgrades that are independent of the supplied
report's main constructions.

The fixed-group observation turns Corvaja--Zannier into a moving-direction
tool.  That uniformity makes it possible to isolate almost the entire
prime-order cyclotomic factor.  The factor then survives reduction of the
whole maximal-contact stencil, not merely one unreduced summand.  Exact
moment-generating and central-difference identities expose the analytic side
of the same construction, and the nearest poles give a variable-order product
formula.  Baker's logarithmic separation shows that the real contact budget
is only logarithmic in the direction height, whereas the surviving
cyclotomic denominator budget is linear in that height.  Consequently, a
large and natural class of growing linear extrapolants is rigorously too weak.
Finally, the adelic covolume calculation shows why simply imposing local
denominator depth cannot restore the missing exponent.

The proof search should now concentrate on mechanisms that are genuinely
nonlinear, tensorial, or multi-gap.  Repeating fixed-order interpolation,
adding more integer nodes, or forcing prescribed denominator divisibility
without a new correlation cannot close the conjecture.

\appendix

\section{A local valuation template for future filters}
\label{app:valuation-template}

The proof of \cref{prop:coefficient-tax} is useful beyond the present
stencil.  Let $K$ be a nonarchimedean field with valuation $v$, and suppose
$b_0,\ldots,b_m$ satisfy
\[
 v(b_m)=-d<0,\qquad v(b_i)\ge0\quad(i<m).
\]
For nonzero coefficients $a_i$ and $S=\sum a_ib_i$, put
$C=\max_i|v(a_i)|$ and $E=\max(0,-v(S))$.  Then $E+C\ge d/2$.
This statement is completely independent of cyclotomic polynomials.  Any
future construction that identifies a prime at which one level has a unique
negative valuation immediately inherits the same coefficient-or-denominator
tax.

\section{Bibliographic notes}

The use of Corvaja--Zannier in \cref{thm:uniform-gcd} is qualitative but
uniform in the fixed finitely generated group; no effective exceptional
height is claimed.  For \eqref{eq:baker-recurrence}, any standard effective
lower bound for a nonzero linear form in two logarithms suffices; the precise
constant is irrelevant.  The period discussion is deliberately diagnostic:
it records the categorical level at which the known theorem stops and does
not assert a new motivic independence result.

\begin{thebibliography}{99}

\bibitem{AnthropicFable2026}
Anthropic,
\emph{Claude Fable 5}, product announcement and system information,
June 9, 2026.
\url{https://www.anthropic.com/claude/fable}.

\bibitem{BakerWustholz1993}
A. Baker and G. W\"ustholz,
\emph{Logarithmic forms and group varieties},
J. Reine Angew. Math. \textbf{442} (1993), 19--62.
\url{https://doi.org/10.1515/crll.1993.442.19}.

\bibitem{ChangPapanikolas2008}
C.-Y. Chang and M. A. Papanikolas,
\emph{Algebraic relations among periods and logarithms of rank $2$ Drinfeld modules},
Amer. J. Math. \textbf{133} (2011), 359--391;
\url{https://arxiv.org/abs/0807.3157}.

\bibitem{CorvajaZannier2005}
P. Corvaja and U. Zannier,
\emph{A lower bound for the height of a rational function at $S$-unit points},
Monatsh. Math. \textbf{144} (2005), 203--224;
\url{https://arxiv.org/abs/math/0311030}.

\bibitem{EpochHadamard2026}
Epoch AI,
\emph{Hadamard matrix of order 668}, FrontierMath open-problem record,
accessed August 22, 2026.
\url{https://epoch.ai/frontiermath/open-problems/hadamard}.

\bibitem{HuberKalck2025}
A. Huber and M. Kalck,
\emph{Dimension formulas for period spaces via motives and species},
arXiv:2405.21053v2, March 2025.
\url{https://arxiv.org/abs/2405.21053}.

\bibitem{HuberWustholz2022}
A. Huber and G. W\"ustholz,
\emph{Transcendence and Linear Relations of 1-Periods},
Cambridge Tracts in Mathematics 227, Cambridge University Press, 2022.

\bibitem{ICARMRank2026}
Institute for Computer-Aided Reasoning in Mathematics,
\emph{Elliptic Curve Rank Leaderboard, curve 273},
record with thirty independent rational points, accessed August 22, 2026.
\url{https://elliptic-rank.icarm.cloud/curve/273}.

\bibitem{Matveev2000}
E. M. Matveev,
\emph{An explicit lower bound for a homogeneous rational linear form in the logarithms of algebraic numbers. II},
Izv. Math. \textbf{64} (2000), 1217--1269.
\url{https://doi.org/10.1070/IM2000v064n06ABEH000314}.

\bibitem{TaoJacobian2026}
T. Tao,
\emph{A digestion of the Jacobian conjecture counterexample},
What's New, July 21, 2026.
\url{https://terrytao.wordpress.com/2026/07/21/a-digestion-of-the-jacobian-conjecture-counterexample/}.

\end{thebibliography}

\end{document}
